% !PRINTED 235
\weberpage{235}
%<<CONT
from which, for $\psi(n)$, the expression
\[
\psi(n) = n\,\Pi\Bigl(1+\frac{1}{p}\Bigr)
\]
results, as in §\,63\footnote{Cf. \gesperrt{Dedekind}: Über die
elliptischen Modulfunktionen. Journal f. Mathematik, Bd.\ 83.}.

By (7) we must now put in the formulae (6) for $n\omega'$:
\[
\frac{\displaystyle \gamma' + \delta'\frac{c+\partial\omega}{a}}{\displaystyle \alpha' + \beta'\frac{c+\partial\omega}{a}},
\]
and the linear transformations of the $f$-functions
[§\,40, (4), (8), (11)] can be applied. In doing so we take

\weq{12}{c \equiv 0\ (\mathrm{mod}\ 16),}

so that, by (3) and (8):
\[
\begin{pmatrix} \alpha', & \beta' \\ \gamma', & \delta' \end{pmatrix} \equiv \begin{pmatrix} a^3, & 0 \\ 0, & a \end{pmatrix}\ (\mathrm{mod}\ 16).
\]

If we denote by $\varrho$ the third root of unity

\weq{13}{\varrho = e^{-\frac{2\pi i}{3}[\alpha'(\gamma'-\beta')-(\alpha'^2-1)\beta'\delta']}\, e^{\frac{2n\pi i}{3}[\alpha(\gamma-\beta)-(\alpha^2-1)\beta\delta]},}

then we obtain
\weq{14}{\begin{aligned}
P_{00} &= \varrho\,\frac{f\left(\dfrac{c+\partial\omega}{a}\right)}{f(\omega)^n},\\
P_{01} &= \varrho\left(\frac{2}{a}\right)\frac{f_1\!\left(\dfrac{c+\partial\omega}{a}\right)}{f_1(\omega)^n},\\
P_{10} &= \varrho\left(\frac{2}{\partial}\right)\frac{f_2\!\left(\dfrac{c+\partial\omega}{a}\right)}{f_2(\omega)^n}.
\end{aligned}}

The quotient

\weq{15}{\frac{P_{10}}{P_{00}} = \left(\frac{2}{\partial}\right)\frac{f_2\left(\dfrac{c+\partial\omega}{a}\right)}{f\left(\dfrac{c+\partial\omega}{a}\right)}\left(\frac{f\omega}{f_2\omega}\right)^n}

gives, by §\,54, (3), if one puts
\[
u(\omega) = \sqrt[4]{k}
\]
the quantities
%CONT>>
